\documentclass[11pt]{article}
\usepackage{amsmath,amssymb,amsthm,hyperref}
\newtheorem{theorem}{Theorem}
\title{A Gaussian adjoint criterion from the whole-source trace bridge}
\author{Internal cross-programme deduction; independent review pending}
\date{15 November 2026}
\begin{document}\maketitle
This draft combines the Gaussian integration-by-parts argument in
\texttt{GAUSSIAN\_KELLER\_ENDOMORPHISMS\_AND\_ONE\_WAY\_TRANSFER.md}
with the whole-source differential-trace argument in
\texttt{WHOLE\_SOURCE\_POWER\_TRACES\_AND\_DIFFERENTIAL\_SPLITTING.md}.
It is not registered as a theorem in the status ledger and makes no novelty
or independent-verification claim. It uses neither Tang classification nor
an assumed counterexample map.

Let $\mathcal L$ be the complex bilinear standard Gaussian Wick functional
on $R=\mathbb C[x_1,\ldots,x_n]$, and put $S(x)=\sum x_i^2/2$.
For $\Phi:\mathbb A^n\to\mathbb A^n$, write $Tf=f\circ\Phi$.

\begin{theorem}
Suppose $\det J\Phi$ is a nonzero constant and $S\circ\Phi-S$ is constant.
Then the following are equivalent:
\begin{enumerate}
\item $\Phi$ is a polynomial automorphism.
\item For every $q\in R$ there is a polynomial $Aq\in\mathbb C[y]$ such that
\[
 \mathcal L_x(q\,Tf)=\mathcal L_y((Aq)f)\quad\text{for every }f\in\mathbb C[y].
\]
\end{enumerate}
If these hold, $Aq=q\circ\Phi^{-1}$.
\end{theorem}
\begin{proof}
Use the polynomial derivations $D_i$ dual to the coordinates of $\Phi$:
$D_i(\Phi_j)=\delta_{ij}$. They exist because the Jacobian determinant is
a scalar unit. Piola's identity gives $\operatorname{div}D_i=0$ and
quadratic preservation gives $D_iS=\Phi_i$. Gaussian integration by parts
therefore gives
\[
 \mathcal L_x(aD_ib)=\mathcal L_x((\Phi_i a-D_i a)b).
 \tag{1}
\]
With $a=1,b=Tf$, these are precisely the normalized Gaussian moment
recurrences for $f\mapsto\mathcal L_x(Tf)$. Their polynomial uniqueness
implies $\mathcal L_x(Tf)=\mathcal L_y(f)$. Thus $T$ preserves the Gaussian
pairing on its image. The pairing on polynomials is nondegenerate, as is
seen in the Hermite basis.

Assume (2). Nondegeneracy makes $A$ unique and complex linear, with
$A1=1$. Moving a target factor through the defining pairing shows
\[
 A((Tg)q)=gAq.\tag{2}
\]
Applying (1) and then target Gaussian integration by parts shows
\begin{align*}
 \mathcal L_y(f A(D_iq))
 &=\mathcal L_x(Tf D_iq)\\
 &=\mathcal L_x(T(y_i f-\partial_i f)q)\\
 &=\mathcal L_y((y_i f-\partial_i f)Aq)
 =\mathcal L_y(f\partial_i(Aq)).
\end{align*}
Consequently $A(D_iq)=\partial_i(Aq)$: $A$ is a differential,
$\mathbb C[y]$-linear retraction of the target inclusion.

For completeness, localize (2) at $F=\mathbb C(y)$. The source becomes
its finite separable function field $L=\mathbb C(x)$.
Nondegeneracy of the field trace pairing writes
$A_F(r)=\operatorname{Tr}_{L/F}(ar)$ for a unique $a\in L$.
Commutation with $D_i$ and field trace gives
$\operatorname{Tr}_{L/F}((D_i a)r)=0$ for every $r$, so $D_i a=0$.
The $D_i$ span the source coordinate derivations; hence $a\in\mathbb C$.
As $A_F(1)=1$, $a=1/[L:F]$. Thus the normalized field trace is polynomial
on the \emph{whole} source ring. In particular all power traces of every
source coordinate are polynomial. Newton identities make those coordinates
integral over the target, so $\Phi$ is finite. A finite complex Keller map
is finite etale and the affine target is simply connected, giving degree
one and a polynomial inverse. This is the retained whole-source trace
argument, not regularity of trace only on a finite normalization.

Conversely, if $\Phi$ has a polynomial inverse, moment preservation applied
to $f(q\circ\Phi^{-1})$ proves the adjoint identity with the asserted $A$.
\end{proof}

Writing $H_\alpha$ for the target Hermite basis, every source $q$ defines
a formal Hermite adjoint with coefficients
\[
 a_\alpha(q)=\frac{\mathcal L_x(q\,H_\alpha(\Phi(x)))}{\alpha!}.
\]
It is polynomial precisely when only finitely many coefficients are
nonzero. Hence any noninvertible map satisfying the theorem's hypotheses
has a fixed polynomial source marker $q$ coupling to arbitrarily high
Hermite degrees of the target. Agreement of every unmarked moment does
not supply polynomial marker transport.

This does \emph{not} exhibit one fixed pure phase $P$ with infinitely many
nonzero $\mathcal L(q(P\circ\Phi)^m)$: the $H_\alpha$ are not powers of
one fixed phase. The endomorphism safety question remains open.
Nor is polynomial adjoint closure claimed for general moment-preserving
maps without the Jacobian and quadratic assumptions.

A separate prospective HC4 connection is the retained claim LCDS1,
\texttt{archive/non-elliptic/plane-jc/LOG\_CONDUCTOR\_DEGREE\_SHIFT.md}.
It gives actual Jacobian-matrix examples with identical determinant and
generic branch Smith profiles but different nodal cokernel gluing.
This is a warning for the current finite-point trace-defect analysis:
identify its actual module before trying to control it by branch data.
No comparison between that conductor module and the new trace-image
module has been established in this draft.
\end{document}
